\scanpage{64}
$a(\gamma(z))$ soit un isom, il f et s que $u^{(1)}$ le soit, i.e. que $u$ soit une isogénie. Alors $u$ admet un inverse $u^{-1}$ unique dans $\operatorname{Hom}_Q(A_Y,B_X)$, \deleted{défini} donc l'hom inverse de $u^{(1)}$ est défini par un $Q$-cycle algébrique. \inserted{$R$, le meilleur dénominateur possible \inserted{étant donné par $\operatorname{card}\operatorname{Ker}u$}}\deleted{\unsure{coeff. en|à coeff. en}} Donc \unsure{$(a(\gamma(z))^{-1})$|$(a(\gamma(z))^{-1}$|$a(\gamma(Z))^{-1}$} est défini par un $Q$-cycle algébrique, ce qui prouve la première assertion, \deleted{La seconde, se voit comme}\inserted{ayant $\unsure{\overline{\mathrm m}|=}$} dénominateur que $R$. Cela prouve la première assertion, la seconde résulte de la connaissance de l'effet cohomologique d'une classe de \inserted{corr.} \unsure{divisoralles|divisorielles}, via l'effet sur les VA \ldots{} (cf \ldots{} ). Il reste à \deleted{prouver}\inserted{noter} que si $D$ admet un certain dénominateur, ce dernier vaut pour le \deleted{hom} $Q$-hom $A_Y\unsure{\mathbin{-}|\to}B_X$ inverse de $u$, donc il \deleted{\unsure{ne cont}} ne peut être meilleur qu'une puissance de \deleted{card} $N$.

\uline{Corollaire} 2.5. Avec les notations de 2.4., pour
% archive sheet number: 63

